% Appendix F: API Data Schemas
\chapter{API DATA SCHEMAS}

This appendix provides the formal Pydantic (Python) models and TypeScript (React) interfaces used for data validation across the platform.

\section{Generic Resource Schema (Python)}
\begin{lstlisting}[language=Python]
class ResourceBase(BaseModel):
    name: str = Field(..., max_length=100)
    provider: ProviderEnum
    resource_type: str
    region: Optional[str] = "us-east-1"
    variables: dict = {}

class ResourceCreate(ResourceBase):
    pass

class ResourceSchema(ResourceBase):
    id: int
    status: str
    provider_id: Optional[str]
    created_at: datetime

    class Config:
        from_attributes = True
\end{lstlisting}

\section{Security Schema (TypeScript)}
\begin{lstlisting}[language=SQL]
interface AuthToken {
    access_token: string;
    token_type: string;
}

interface UserProfile {
    id: number;
    username: string;
    email: string;
    is_active: boolean;
}
\end{lstlisting}

\newpage
% Appendix G: Project Glossary
\chapter{PROJECT GLOSSARY}

\section{Cloud Terms}
\begin{itemize}
    \item \textbf{AWS (Amazon Web Services)}: A subsidiary of Amazon providing on-demand cloud computing platforms and APIs.
    \item \textbf{Azure}: A cloud computing service created by Microsoft for building, testing, deploying, and managing applications.
    \item \textbf{GCP (Google Cloud Platform)}: A suite of cloud computing services that runs on the same infrastructure that Google uses internally.
    \item \textbf{Multi-Cloud}: The use of multiple cloud computing and storage services in a single heterogeneous architecture.
    \item \textbf{Vendor Lock-in}: A situation in which a customer is dependent on a vendor for products and services and cannot switch to another vendor without substantial costs.
\end{itemize}

\section{Technical Terms}
\begin{itemize}
    \item \textbf{FastAPI}: A modern, high-performance web framework for building APIs with Python 3.7+ based on standard Python type hints.
    \item \textbf{Terraform}: An open-source infrastructure as code software tool created by HashiCorp.
    \item \textbf{Celery}: An asynchronous task queue/job queue based on distributed message passing.
    \item \textbf{Redis}: An in-memory data structure store used as a distributed, in-memory key-value database, cache, and message broker.
    \item \textbf{JWT (JSON Web Token)}: An open standard that defines a compact and self-contained way for securely transmitting information between parties as a JSON object.
    \item \textbf{AES-128}: Advanced Encryption Standard with a 128-bit key length, used for securing cloud credentials in the Nebula vault.
\end{itemize}

\newpage
